% Part: proof-theory
% Chapter: natural-deduction
% Section: translation-N2i

\documentclass[../../../include/open-logic-section]{subfiles}

\begin{document}

\olfileid{pt}{nat}{tni}

\olsection{الترجمة من \Log{N2} إلى \Log{G2}}

\begin{prop}\ollabel{prop:N2-to-G2}
إذا كان $\Log{N2c}\ (\Log{N2i}) \Proves \Gamma \Sequent !A$، فإن
$\Log{G2c}\ (\Log{G2i}) + \Cut \Proves \Gamma' \Sequent \Delta$، حيث
$\Gamma'$ هي المجموعة المتعددة من~!!{formula}s الناتجة من~$\Gamma$ بحذف
الوسوم، وحيث $\Delta = \{!A\}$ إذا لم تكن~$!A$ هي~$\lfalse$، و
$\Delta = \emptyset$ إذا كانت إياها.
\end{prop}

\begin{proof}
البرهان مرة أخرى بالاستقراء في ارتفاع~!!{proof}~$\delta$ للمتتالية
$\Gamma \Sequent !A$. فإذا كانت~$\delta$ مسلّمة
$x:!A \Sequent !A$، كان~$\pi$ هو المسلّمة المقابلة
$!A \Sequent !A$، أو~$\lfalse \Sequent \ $ إذا كان
$!A \ident \lfalse$.

إذا انتهى~$\delta$ بقاعدة استدلال، نميّز الحالات الآتية:
\begin{enumerate}
  \item إذا كانت~$R$ هي~\Intro{\lif} وصيغتها الرئيسة~$!B \lif !C$، وكان
  $x:!B$ موجودًا في سياق المقدمة دون النتيجة، انتهى~$\delta$ بما يأتي:
  \[
  \AxiomC{}
  \RightLabel{$\delta_1$}
  \Deduce$x:!B, \Gamma \fCenter !C$
  \RightLabel{\Intro{\lif}}
  \UnaryInf$\Gamma \fCenter !B \lif !C$
  \DisplayProof
  \]
  يعطي فرض الاستقراء~!!a{proof}~$\pi_1$ للمتتالية
  $!B,\Gamma' \Sequent !C$، أو للمتتالية~$!B,\Gamma' \Sequent \ $ إذا
  كان $!C \ident \lfalse$. وفي الحالة الأخيرة نطبّق أولًا
  \RightR{\Weakening} لإضافة~$\lfalse$ إلى الطرف الأيمن. ثم نحصل على
  $\pi$ في الحالتين بتطبيق~\RightR{\lif}.
  \item إذا كانت~$R$ هي~\Intro{\lif} لكن~$x:!B$ لم يكن موجودًا في سياق
  المقدمة، انتهى~$\delta$ بما يأتي:
  \[
  \AxiomC{}
  \RightLabel{$\delta_1$}
  \Deduce$\Gamma \fCenter !C$
  \RightLabel{\Intro{\lif}}
  \UnaryInf$\Gamma \fCenter !B \lif !C$
  \DisplayProof
  \]
  يعطي فرض الاستقراء~!!a{proof}~$\pi_1$ للمتتالية
  $\Gamma' \Sequent !C$، أو للمتتالية ذات الطرف الأيمن الخالي إذا كان
  $!C \ident \lfalse$. في الحالة الأخيرة نطبّق أولًا
  \RightR{\Weakening} لإضافة~$\lfalse$، ثم نطبّق
  \LeftR{\Weakening} بالصيغة~$!B$، وأخيرًا~\RightR{\lif}؛ وفي الحالة
  الأولى نحتاج إلى التطبيقين الأخيرين فقط.
  \item إذا كانت~$R$ هي~\Elim{\lif}، انتهى~$\delta$ بما يأتي:
  \[
  \AxiomC{}
  \RightLabel{$\delta_1$}
  \Deduce$\Gamma_1 \fCenter !B \lif !A$
  \AxiomC{}
  \RightLabel{$\delta_2$}
  \Deduce$\Gamma_2 \fCenter !B$
  \RightLabel{\Elim{\lif}}
  \BinaryInf$\Gamma_1, \Gamma_2 \fCenter !A$
  \DisplayProof
  \]
  حيث السياق~$\Gamma$ هو تجميع السياقين~$\Gamma_1,\Gamma_2$. يعطي فرض
  الاستقراء !!{proof}s: برهانًا~$\pi_1$ للمتتالية
  $\Gamma_1' \Sequent !B \lif !A$، وبرهانًا~$\pi_2$ للمتتالية
  $\Gamma_2' \Sequent !B$ إذا كان~$!B \not\ident \lfalse$، أو للمتتالية
  $\Gamma_2' \Sequent \ $ إذا كان~$!B \ident \lfalse$. وفي الحالة
  الثانية نضيف~\RightR{\Weakening} إلى~$\pi_2$ كي تكون نتيجته
  $\Gamma_2' \Sequent \lfalse$.

  إذا كان $!A \not\ident \lfalse$ نحصل على~!!{proof}~$\pi$ بالصورة:
  \[
  \AxiomC{}
  \RightLabel{$\pi_1$}
  \Deduce$\Gamma_1' \fCenter !B \lif !A$
  \AxiomC{}
  \RightLabel{$\pi_2$}
  \Deduce$\Gamma_2' \fCenter !B$
  \Axiom$!A \fCenter !A$
  \RightLabel{\LeftR{\lif}}
  \BinaryInf$!B \lif !A, \Gamma_2' \fCenter !A$
  \RightLabel{\Cut}
  \BinaryInf$\Gamma_1', \Gamma_2' \fCenter !A$
  \DisplayProof
  \]
  أما إذا كان $!A \ident \lfalse$ فنستخدم~$\lfalse \Sequent \ $ مقدمةً
  يمنى لتطبيق~\LeftR{\lif}، فنحصل على~!!a{proof} للطرف الأيمن الخالي، ثم
  نطبّق~\Cut؛ وهكذا يكون~!!{proof}~$\pi$ للمتتالية
  $\Gamma_1',\Gamma_2' \Sequent \ $ ولا نطبّق إضعافًا يمينيًا عليها.
  قد لا يساوي تجميع المجموعتين المتعددتين~$\Gamma_1',\Gamma_2'$ المجموعة
  المتعددة~$\Gamma'$ الناتجة من تجميع~$\Gamma_1,\Gamma_2$. فمثلًا، إذا كان
  كل من~$\Gamma_1$ و$\Gamma_2$ يحتوي~$x:!C$، احتوى التجميع الأول نسختين من
  $!C$، بينما لا ينتج من تجميع السياقين الموسومين إلا نسخة واحدة. نصحح ذلك
  بإضافة استدلالات مناسبة من~\LeftR{\Contraction}.

\item إذا كانت~$R$ هي~\FalseCl، انتهى~$\delta$ بما يأتي:
  \[
  \AxiomC{}
  \RightLabel{$\delta_1$}
  \Deduce$x: \lnot !A, \Gamma \fCenter \lfalse$
  \RightLabel{\FalseCl}
  \UnaryInf$\Gamma \fCenter !A$
  \DisplayProof
  \]
  يوجد، بفرض الاستقراء، !!a{proof}~$\pi_1$ للمتتالية
  $\lnot !A,\Gamma' \Sequent \ $. ونتخذ~!!{proof}~$\pi$ ليكون:
  \[
  \Axiom$!A \fCenter !A$
  \RightLabel{\RightR{\lnot}}
  \UnaryInf$\fCenter !A, \lnot !A$
  \AxiomC{}
  \RightLabel{$\pi_1$}
  \Deduce$\lnot !A, \Gamma' \fCenter $
  \RightLabel{\Cut}
  \BinaryInf$\Gamma' \fCenter !A$
  \DisplayProof
  \]
  \item تُعالج قواعد الاستدلال المتبقية على نحو مماثل، مع تجميع السياقات
  بالفواصل وحذف التكرارات الزائدة بانقباضات يسارية عند الحاجة.
\end{enumerate}
لاحظ أن ترجمة~\FalseCl{} وحدها تنتج~!!a{proof}~$\pi$ فيه أكثر من
!!{formula} واحدة على اليمين. أي إن ترجمة~!!{proof} في~\Log{N2i} هي
!!{proof} في~\Log{G2i}.
\end{proof}

\end{document}
